\documentclass{article}
\usepackage[top=2cm,bottom=2cm,left=3cm,right=3cm]{geometry}
\usepackage{enumitem}
\usepackage{xcolor}
\usepackage{minted}
\usepackage[most]{tcolorbox}
\tcbuselibrary{minted,breakable}
\usepackage{fontspec}
\usepackage{polyglossia}
\setdefaultlanguage{ukrainian}
\setotherlanguages{english}
\usepackage{Alegreya}
\usepackage{FiraMono}
\newfontfamily\cyrillicfont{Alegreya}
\newfontfamily\cyrillicfontsf{Alegreya Sans}
\newfontfamily\cyrillicfonttt{DejaVu Sans Mono}
\usepackage{sciffi-python}
\usepackage{sciffi-memo}

\usemintedstyle{emacs}

\newtcolorbox{CodeBox}[1][]{
    breakable,
    enhanced,
    colback=gray!10,
    colframe=gray!70,
    fonttitle=\bfseries,
    title=#1,
    sharp corners,
    boxrule=0.8pt,
    top=2mm,
    bottom=2mm,
    left=2mm,
    right=2mm,
    fontupper=\footnotesize,
}

\newtcolorbox{ResultBox}[1][]{
    breakable,
    enhanced,
    colback=gray!5,
    colframe=gray!50,
    fonttitle=\bfseries,
    title=#1,
    sharp corners,
    boxrule=0.6pt,
    top=1mm,
    bottom=1mm,
    left=2mm,
    right=2mm,
    fontupper=\footnotesize,
}

\begin{document}
    \begin{titlepage}
        \begin{center}
            {\Large
                МІНІСТЕРСТВО ОСВІТИ ТА НАУКИ УКРАЇНИ

                НАЦІОНАЛЬНИЙ ТЕХНІЧНИЙ УНІВЕРСИТЕТ УКРАЇНИ
                «КИЇВСЬКИЙ ПОЛІТЕХНІЧНИЙ ІНСТИТУТ ІМЕНІ ІГОРЯ СІКОРСЬКОГО»

                ФАКУЛЬТЕТ ІНФОРМАТИКИ ТА ОБЧИСЛЮВАЛЬНОЇ ТЕХНІКИ

                КАФЕДРА ІНФОРМАЦІЙНИХ СИСТЕМ ТА ТЕХНОЛОГІЙ
            }


            \vspace{60mm}
            {\large
                \textbf{ЗВІТ}

                \vspace{5mm}

                \textbf{До комп'ютерного практикуму №5}

                \vspace{5mm}

                з дисципліни «Технології Розробки Програмного Забезпечення»
            }

            На тему «Операції об’єднання результатів декількох запитів. Модифікація даних.»

            варіант \textnumero 14 (\textnumero 4)
        \end{center}

        \vfill
        \hfill
        \begin{minipage}[t]{0.35\textwidth}
            \textbf{Виконав:}

            Носилевський В. А., ІК-32
        \end{minipage}

        \vfill

        \begin{center}
            {\bf
                Київ

                КПІ Ім. Ігоря Сікорського

                2025
            }
        \end{center}
    \end{titlepage}
    \newpage

    \textbf{Тема:} операції об’єднання результатів декількох запитів. Модифікація даних.

    \textbf{Мета:} вивчити операції об’єднання результатів декількох запитів в одну таблицю (UNION, INTERSECT, EXCEPT) та команди модифікації даних (INSERT, UPDATE, DELETE).

    \textbf{Навички що будуть здобуті:} вміння об’єднувати запити до бази даних, використовувати копіювання, оновлення та видалення даних.
    
    \section*{Хід виконання роботи}
    \begin{sciffi}{python}
        from pathlib import Path
        from subprocess import run, PIPE
        from collections.abc import Sequence
        from contextlib import contextmanager
        
        @contextmanager
        def env(name: str, *args: str):
            print(f"\\begin{{{name}}}" + "".join(args))
            try:
                yield
            finally:
                print(f"\\end{{{name}}}")

        def sh(*args: str | Path) -> str:
            return run(
                tuple(map(str, args)),
                stdout=PIPE,
                stderr=PIPE,
                text=True,
            ).stdout

        scripts = Path("./scripts/")

        flow: Sequence[tuple[Path, Sequence[str]]] = [
            (script, ("--no-echo",))
            for script in sorted(
                scripts.glob("*.sql"),
                key=lambda p: int(p.stem),
            )
        ]

        for script, args in flow:
            source = script.read_text(encoding="utf-8")
            output = sh("tsqlexec", *args, script)

            print(f"\\subsection*{{Запит з файлу: {script.name}}}")
            with env("CodeBox", "[SQL Source]"):
                with env("minted", "[breaklines,linenos]", "{sql}"):
                    print(source, end="")
            print()

            if not output.strip():
                continue

            print("Result:")
            with env("ResultBox", "[Output]"):
                with env("minted", "[breaklines]", "{text}"):
                    print(output, end="")
            print("\\vspace{5mm}")
    \end{sciffi}

    \newpage
    \section*{Контрольні питання}
    \begin{enumerate}[label=\arabic*., leftmargin=*]
        \item \textbf{Які вимоги пред’являються до таблиць, над якими виконуються оператори UNION?}
        
        Для успішного виконання оператора \texttt{UNION} (а також \texttt{INTERSECT} та \texttt{EXCEPT}) всі запити \texttt{SELECT}, що об'єднуються, повинні відповідати таким вимогам:
        \begin{itemize}
            \item \textbf{Кількість стовпців:} Усі запити повинні повертати однакову кількість стовпців.
            \item \textbf{Сумісність типів даних:} Типи даних відповідних стовпців (за їх порядком) у кожному запиті мають бути сумісними. Це не означає, що вони мають бути ідентичними, але СКБД повинна мати можливість неявно перетворити їх до одного спільного типу.
        \end{itemize}
        Назви стовпців у результуючій таблиці беруться з першого запиту \texttt{SELECT}.

        \item \textbf{Поняття та синтаксис UNION?}
        
        Оператор \texttt{UNION} використовується для об'єднання результатів двох або більше запитів \texttt{SELECT} в один результуючий набір. За замовчуванням, \texttt{UNION} видаляє всі дублікати рядків з кінцевого результату.
        \begin{CodeBox}[Синтаксис UNION]
        \begin{minted}[breaklines]{sql}
SELECT column1, column2 FROM table1
UNION
SELECT columnA, columnB FROM table2;
        \end{minted}
        \end{CodeBox}
        
        \item \textbf{В чому різниця між UNION та UNION ALL?}
        
        Основна різниця полягає у роботі з дублікатами:
        \begin{itemize}
            \item \textbf{UNION:} Видаляє дублікати. Для цього СКБД повинна виконати додаткову роботу (сортування або хешування), що робить цю операцію повільнішою.
            \item \textbf{UNION ALL:} Включає всі рядки з усіх запитів, включаючи дублікати. Ця операція є значно швидшою, оскільки не потребує перевірки на унікальність.
        \end{itemize}

        \item \textbf{Поняття та синтаксис INTERSECT?}
        
        Оператор \texttt{INTERSECT} повертає тільки ті унікальні рядки, які є спільними для обох запитів \texttt{SELECT} (тобто присутні в результатах і першого, і другого запиту).
        \begin{CodeBox}[Синтаксис INTERSECT]
        \begin{minted}[breaklines]{sql}
SELECT column_name(s) FROM table1
INTERSECT
SELECT column_name(s) FROM table2;
        \end{minted}
        \end{CodeBox}

        \item \textbf{Поняття та синтаксис EXCEPT?}
        
        Оператор \texttt{EXCEPT} повертає унікальні рядки з результату першого запиту, яких немає в результаті другого запиту. Він виконує "віднімання" одного набору результатів від іншого.
        \begin{CodeBox}[Синтаксис EXCEPT]
        \begin{minted}[breaklines]{sql}
SELECT column_name(s) FROM table1
EXCEPT
SELECT column_name(s) FROM table2;
        \end{minted}
        \end{CodeBox}
        
        \item \textbf{Поняття та синтаксис INSERT?}
        
        Команда \texttt{INSERT INTO} використовується для додавання одного або кількох нових рядків до таблиці.
        \begin{CodeBox}[Синтаксис INSERT]
        \begin{minted}[breaklines]{sql}
-- Варіант 1: Вказуються стовпці (рекомендовано)
INSERT INTO table_name (column1, column2) 
VALUES (value1, value2);

-- Варіант 2: Стовпці не вказуються (потрібно надати значення для всіх полів у правильному порядку)
INSERT INTO table_name 
VALUES (value1, value2, value3, ...);
        \end{minted}
        \end{CodeBox}

        \item \textbf{INSERT INTO для таблиць з однаковою структурою в одній БД?}
        
        Можна скопіювати всі дані з однієї таблиці в іншу за допомогою конструкції \texttt{INSERT INTO ... SELECT}.
        \begin{CodeBox}[Копіювання даних]
        \begin{minted}[breaklines]{sql}
INSERT INTO TargetTable
SELECT * FROM SourceTable
WHERE condition; -- Умова WHERE не є обов'язковою
        \end{minted}
        \end{CodeBox}

        \item \textbf{INSERT INTO для таблиць з різною структурою в одній БД?}
        
        Якщо структури відрізняються, потрібно явно вказати відповідність стовпців.
        \begin{CodeBox}[Копіювання з різною структурою]
        \begin{minted}[breaklines]{sql}
INSERT INTO TargetTable (TargetColumn1, TargetColumn2)
SELECT SourceColumnA, SourceColumnB 
FROM SourceTable;
        \end{minted}
        \end{CodeBox}

        \item \textbf{INSERT INTO для таблиць з однаковою структурою в різних БД?}
        
        Для звернення до таблиці в іншій базі даних використовується повне ім'я, що складається з трьох частин: \texttt{database.schema.table}.
        \begin{CodeBox}[Копіювання між БД]
        \begin{minted}[breaklines]{sql}
INSERT INTO TargetDB.dbo.TargetTable
SELECT * FROM SourceDB.dbo.SourceTable;
        \end{minted}
        \end{CodeBox}
        
        \item \textbf{SELECT <...> INTO для таблиць з однаковою структурою в одній БД?}
        
        Конструкція \texttt{SELECT ... INTO} створює \textbf{нову} таблицю і копіює в неї результат запиту. Цільова таблиця не повинна існувати на момент виконання команди.
        \begin{CodeBox}[Створення нової таблиці з даними]
        \begin{minted}[breaklines]{sql}
SELECT *
INTO NewTable
FROM ExistingTable
WHERE condition;
        \end{minted}
        \end{CodeBox}

        \item \textbf{SELECT <...> INTO для таблиць з однаковою структурою в різних БД?}
        
        Аналогічно до \texttt{INSERT}, для створення таблиці в іншій БД використовується повне ім'я.
        \begin{CodeBox}[Створення нової таблиці в іншій БД]
        \begin{minted}[breaklines]{sql}
SELECT *
INTO TargetDB.dbo.NewTable
FROM SourceDB.dbo.ExistingTable;
        \end{minted}
        \end{CodeBox}
        
        \item \textbf{Поняття та синтаксис UPDATE?}
        
        Команда \texttt{UPDATE} використовується для зміни існуючих записів у таблиці.
        \begin{CodeBox}[Синтаксис UPDATE]
        \begin{minted}[breaklines]{sql}
UPDATE table_name
SET column1 = value1, column2 = value2
WHERE condition;
        \end{minted}
        \end{CodeBox}
        Ключове слово \texttt{WHERE} є критично важливим, оскільки без нього команда оновить \textbf{усі} рядки в таблиці.
        
        \item \textbf{Як виконати UPDATE, якщо умова змін відноситься до пов’язаної таблиці?}
        
        Існує два основних способи: використання підзапиту або \texttt{JOIN}.
        \begin{CodeBox}[UPDATE з підзапитом]
        \begin{minted}[breaklines]{sql}
UPDATE Products
SET Price = Price * 1.1
WHERE CategoryID IN (SELECT ID FROM Categories WHERE Name = 'Electronics');
        \end{minted}
        \end{CodeBox}
        \begin{CodeBox}[UPDATE з JOIN (синтаксис T-SQL)]
        \begin{minted}[breaklines]{sql}
UPDATE p
SET p.Price = p.Price * 1.1
FROM Products AS p
JOIN Categories AS c ON p.CategoryID = c.ID
WHERE c.Name = 'Electronics';
        \end{minted}
        \end{CodeBox}

        \item \textbf{Поняття та синтаксис DELETE?}
        
        Команда \texttt{DELETE} використовується для видалення одного або кількох існуючих рядків з таблиці.
        \begin{CodeBox}[Синтаксис DELETE]
        \begin{minted}[breaklines]{sql}
DELETE FROM table_name WHERE condition;
        \end{minted}
        \end{CodeBox}
        Без умови \texttt{WHERE} команда видалить \textbf{усі} рядки з таблиці.
        
        \item \textbf{Як виконати DELETE, якщо умова видалення відноситься до пов’язаної таблиці?}
        
        Так само як і з \texttt{UPDATE}, можна використовувати підзапити або \texttt{JOIN}.
        \begin{CodeBox}[DELETE з підзапитом]
        \begin{minted}[breaklines]{sql}
DELETE FROM Orders
WHERE CustomerID IN (SELECT ID FROM Customers WHERE Country = 'Ukraine');
        \end{minted}
        \end{CodeBox}
        \begin{CodeBox}[DELETE з JOIN (синтаксис T-SQL)]
        \begin{minted}[breaklines]{sql}
DELETE o
FROM Orders AS o
JOIN Customers AS c ON o.CustomerID = c.ID
WHERE c.Country = 'Ukraine';
        \end{minted}
        \end{CodeBox}
        
    \end{enumerate}
\end{document}
